\heading{Part 1: a project you can resume}{Summary \hfill 28 September 2026}
\textbf{Keep the project's memory in files.} Ask the agent to inspect and explain;
check its claims against the files yourself. Your goal is a small, recoverable
checkpoint, not a finished software package.

\topic{The workflow}
\textbf{Inspect $\rightarrow$ save $\rightarrow$ restart $\rightarrow$ compare $\rightarrow$ hand off.}
A \textbf{baseline} is a preserved record of what the original computation
produced. Compare a fresh run against it before changing the model.
Repeatability means the run agrees; it does not establish scientific validity.

\topic{Today's case: a cooled reactor}
The CSTR (continuous stirred-tank reactor) notebook computes steady concentration
and temperature as coolant temperature changes. You do not need to derive the
model today. Locate its assumptions, units, inputs, and outputs before trusting
an agent's explanation. Keep notebook code, parameters, and solver settings unchanged.
A \textbf{nominal} record uses the reference coolant condition; a \textbf{sweep}
varies coolant temperature. A branch index identifies a state at one condition.

\renewcommand{\arraystretch}{1.12}
\begin{tabularx}{\linewidth}{@{}p{2.8in}X@{}}
\toprule
\textbf{In your starter folder} & \textbf{Use it to find}\\
\midrule
\code{notebooks/cstr_exploration.ipynb} & Model, computation, assumptions\\
\code{data/reactor_parameters.yml} & Parameter values and units\\
\code{AGENTS.md} or \code{CLAUDE.md} & Your short project instructions\\
\code{results/baseline.json} & Captured values, settings, provenance\\
\code{results/steady_states.csv} & Sweep records; locate by coolant temperature and branch index\\
\code{results/baseline.executed.ipynb} & Saved executed notebook\\
\code{results/recheck/comparison.json} & Fresh run versus saved baseline\\
\code{docs/handoff.md} & Checked work, pending work, next step\\
\bottomrule
\end{tabularx}

\savework

\vfill
\textbf{Before the meeting:} use the starter README to create and activate its
environment in your own private copy. In the full checkout, the starter is
\code{resources/workshops/cstr/}. All activity commands run from that folder.\par
\small\repourl\par
\textbf{Turn over for the numbered activity.} Keep this sheet beside your keyboard.
